%% kapa: exam=2 round=31 year=2020 period=2 subject=appraisal_theory subject_ko=감정평가이론 seq=3 number=3 points=20
%% source: Q-Net 감정평가사 기출문제 (cst00309) artlSeq=5207673 file=raw/1008/5207673/2교시_감정평가이론.pdf
%% conversion: kordoc 4.24.51 -> tools/gen2.py (auto); review=proofread-gpt-astra; flags=none
\begin{longquestion}{【문제 3】}
A 토지는 ○○재개발사업구역에 소재하고 있다. A 토지에 대하여 재개발사업의 절차상 종전자산의 감정평가를 하는 경우와 손실보상(현금청산)을 위한 감정평가를 하는 경우에 다음의 물음에 답하시오. (20점) \par
\medskip\noindent 물음 1) 각각의 감정평가에 있어 기준시점, 감정평가액의 성격 및 감정평가액 결정 시 고려할 점에 관하여 설명하시오. (10점) \par
\medskip\noindent 물음 2) 각각의 감정평가에 있어 재개발사업으로 인한 개발이익의 반영여부에 관하여 설명하시오. (10점) \par
\end{longquestion}
